\section{Contrôle thêta du réseau de Leech : identité \texorpdfstring{\(\Theta_{\Lambda}=E_4^3-720\Delta\)}{Theta Lambda = E4^3 - 720 Delta} et séparation entre coïncidence arithmétique et morphisme structurel}
\label{sec:control-theta-leech}

À partir de \(\Lambda_{24}\), la chaîne
\[
\Lambda_{24}
\longrightarrow V^\natural
\longrightarrow\mathbb M
\longrightarrow\text{Moonshine}
\]
relève des mathématiques classiques : algèbre d'opérateurs de vertex, groupe Monstre et
fonctions de McKay--Thompson. La contribution vérifiée ici se situe dans l'entrée
finie, la sélection du bord et le \(3\)-voisin sans racines, non dans une nouvelle
appropriation de ces théorèmes.

Trois identités numériques accompagnent cette chaîne :
\[
\Theta_{\Lambda}=E_4^3-720\Delta,
\]
\[
196884=196560+324,
\]
\[
196884=54(3645+1)=54(5\cdot729+1).
\]
Ce sont des identités arithmétiques exactes, mais sans morphisme ni décomposition
de caractères, elles ne constituent pas à elles seules des signes mathématiques de
compatibilité. Aucune ne remplace la construction de \(V^\natural\) ni
ne fournit une action directe du Monstre sur les mots décimaux HMT.
